%%%%%%%%%%%%%%%%%%%%%%%%%%%%%%%%%%%%%%%%%%%%%%%%%%%%%%%%%%%%%%%%%%%%%
%% sec_2_2_EM_interaction.tex
%%
%% Section 2.2 for paper2_main.tex
%% RMG Periodic RT-TDDFT -- Paper 2
%%
%% Covers:
%%   2.2  Interaction of Matter with Electromagnetic Radiation
%%        2.2.1  General TDKS equation with EM coupling
%%        2.2.2  The electromagnetic plane wave
%%        2.2.3  Multipole expansion and long-wavelength approximation
%%        2.2.4  Failure of the length gauge for periodic systems
%%        2.2.5  Scope and limitations of the velocity gauge
%%        2.2.6  Beyond the dipole approximation: the full vector potential
%%        2.2.7  Observables: dipole moment, current density, Berry phase
%%        2.2.8  Summary and scope of the present work
%%
%% CONVENTIONS:
%%   Atomic units: hbar = m_e = e = 4*pi*eps_0 = 1
%%   Electron charge = -1  (e > 0 is elementary charge magnitude)
%%   Imaginary unit: \imath  (dotless i, consistent with paper 1)
%%   \myblue{} for all annotations and placeholders
%%   NEVER use \myred{}
%%
%% COMPILE NOTE:
%%   This file is intended to be \input{} into paper2_main.tex.
%%   It assumes the following macros are defined in the preamble:
%%     \Avec, \Evec, \Bvec, \phat, \VKS, \iunit
%%   If these are not available, uncomment the fallback block below.
%%%%%%%%%%%%%%%%%%%%%%%%%%%%%%%%%%%%%%%%%%%%%%%%%%%%%%%%%%%%%%%%%%%%%

%%--- Fallback macro definitions (uncomment if not defined in preamble) ---
\providecommand{\Avec}{\mathbf{A}}
\providecommand{\Evec}{\mathbf{E}}
\providecommand{\Bvec}{\mathbf{B}}
\providecommand{\phat}{\hat{\mathbf{p}}}
\providecommand{\VKS}{V_{\text{eff}}}
\providecommand{\iunit}{\imath}
\providecommand{\myblue}[1]{\textcolor{blue}{#1}}


%--------------------------------------------------------------------
\subsection{Interaction of Matter with Electromagnetic Radiation}
\label{sec:em-interaction}
%--------------------------------------------------------------------

The starting point for any RT-TDDFT calculation is the coupling of
the electronic system to an external time-dependent perturbation.
This section establishes the general framework for the interaction
of electrons with electromagnetic radiation, motivates the specific
choices of gauge and approximation used in this work, and identifies
the observables through which spectral properties are extracted from
the simulation.


%--------------------------------------------------------------------
\subsubsection{The general time-dependent Kohn-Sham equation with
               electromagnetic coupling}
\label{sec:tdks-em}
%--------------------------------------------------------------------

The starting point for RT-TDDFT in the presence of an electromagnetic
field is the time-dependent Kohn-Sham (TDKS) equation with the
standard minimal coupling prescription.  For a single KS orbital
$\psi$ of an electron (charge $q_e = -e$, with $e > 0$), the TDKS
equation reads
\begin{equation}
  \iunit\hbar\frac{\partial\psi(\mathbf{r},t)}{\partial t}
  = \hat{H}(t)\,\psi(\mathbf{r},t)
  \label{eq:tdks-general}
\end{equation}
with the Hamiltonian
\begin{equation}
  \hat{H}(t)
  = \frac{1}{2m}
    \!\left(\phat - \frac{q_e}{c}\Avec(\mathbf{r},t)\right)^{\!2}
  + q_e\,\Phi(\mathbf{r},t)
  + V_\text{eff}[\rho](\mathbf{r},t).
  \label{eq:H-general}
\end{equation}
Here $\Avec(\mathbf{r},t)$ and $\Phi(\mathbf{r},t)$ are the vector
and scalar potentials describing the electromagnetic field, and
$V_\text{eff}[\rho]$ is the Kohn-Sham effective potential containing
the ionic, Hartree, and exchange-correlation contributions.

In the general form, two electromagnetic potentials ($\Avec$ and
$\Phi$) and the KS potential all appear simultaneously.  However,
the gauge freedom of electrodynamics---the fact that $\Avec$ and
$\Phi$ can be transformed without changing the physical fields
$\Evec$ and $\Bvec$---allows us to simplify significantly.
The choice of gauge determines how the EM coupling enters the
Hamiltonian.

Throughout this work we adopt the \textbf{Coulomb gauge}, defined
by two conditions:
\begin{equation}
  \nabla\cdot\Avec = 0
  \quad\text{(transversality)},
  \qquad
  \Phi_\text{rad} = 0
  \quad\text{(no scalar potential from radiation)}.
  \label{eq:coulomb-gauge}
\end{equation}
In this gauge, all information about the radiation field is carried
exclusively by $\Avec(\mathbf{r},t)$.  The only scalar potential
that remains is the instantaneous Coulomb interaction between
charges, which is already contained in $V_\text{eff}$.
The TDKS Hamiltonian simplifies to
\begin{equation}
  \hat{H}(t)
  = \frac{1}{2m}
    \!\left(\phat + \frac{e}{c}\Avec(\mathbf{r},t)\right)^{\!2}
  + V_\text{eff}[\rho](\mathbf{r},t),
  \label{eq:H-coulomb}
\end{equation}
where the sign flip ($-q_e \to +e$) comes from $q_e = -e$ for the
electron.

Expanding the squared kinetic term gives three physically distinct
contributions:
\begin{equation}
  \hat{H}(t)
  = \underbrace{\frac{\hat{p}^2}{2m}}_{T_0}
  + \underbrace{\frac{e}{mc}\,\Avec\cdot\phat}_{\text{paramagnetic}}
  + \underbrace{\frac{e^2}{2mc^2}\,\Avec^2}_{\text{diamagnetic}}
  + V_\text{eff}.
  \label{eq:H-expanded}
\end{equation}
The \emph{paramagnetic term} is the leading-order light-matter
coupling; it is linear in $\Avec$ and couples the field to the
electronic current.
The \emph{diamagnetic term} is quadratic in $\Avec$ and is typically
negligible for weak perturbative fields (it enters at second order
in perturbation theory), though it is required for gauge invariance
of the response and for satisfaction of the f-sum rule.
In the Coulomb gauge, $\nabla\cdot\Avec = 0$ guarantees that
$\Avec\cdot\phat = \phat\cdot\Avec$, so the paramagnetic term
simplifies to a single $(e/mc)\Avec\cdot\phat$ without
symmetrization.

\myblue{[Molecular analogy:]}
In Gaussian basis set codes, the same Hamiltonian appears but
$\phat$ is represented as matrix elements
$\langle\phi_\mu|\phat|\phi_\nu\rangle$ between atomic orbitals.
In the real-space grid representation used by RMG, $\phat =
-\iunit\hbar\nabla$ acts directly through finite-difference stencils
on the grid---no basis transformation is needed.


%--------------------------------------------------------------------
\subsubsection{The electromagnetic plane wave}
\label{sec:plane-wave}
%--------------------------------------------------------------------

A monochromatic EM plane wave in the Coulomb gauge is described by
the vector potential
\begin{equation}
  \Avec(\mathbf{r},t)
  = A_0\,\hat{\varepsilon}\,
    e^{\iunit(\mathbf{q}\cdot\mathbf{r}-\omega t)}
    + \text{c.c.},
  \label{eq:plane-wave-A}
\end{equation}
where $\hat{\varepsilon}$ is the polarization unit vector,
$\mathbf{q}$ is the photon wavevector with $|\mathbf{q}|=\omega/c$,
and $\omega$ is the angular frequency.
The transversality condition $\nabla\cdot\Avec=0$ requires
$\hat{\varepsilon}\cdot\mathbf{q}=0$: the polarization is
perpendicular to the direction of propagation.

The physical electric and magnetic fields are obtained from $\Avec$
alone through Maxwell's relations:
\begin{align}
  \Evec(\mathbf{r},t)
  &= -\frac{1}{c}\frac{\partial\Avec}{\partial t}
   = \frac{\iunit\omega A_0}{c}\,\hat{\varepsilon}\,
     e^{\iunit(\mathbf{q}\cdot\mathbf{r}-\omega t)}
     + \text{c.c.},
  \label{eq:E-from-A}
  \\[4pt]
  \Bvec(\mathbf{r},t)
  &= \nabla\times\Avec
   = \iunit A_0\,(\mathbf{q}\times\hat{\varepsilon})\,
     e^{\iunit(\mathbf{q}\cdot\mathbf{r}-\omega t)}
     + \text{c.c.}
  \label{eq:B-from-A}
\end{align}
This confirms that a single vector potential in the Coulomb gauge is
sufficient to represent both physical fields---no scalar potential
is needed.  In Gaussian units,
$|\Bvec|/|\Evec| = 1$ for a plane wave, but the coupling of
non-relativistic electrons to $\Bvec$ is weaker than to $\Evec$ by
a factor of $v/c \approx \alpha \approx 1/137$.  This hierarchy is
the physical origin of the dominance of electric-dipole transitions
in ordinary spectroscopy.


%--------------------------------------------------------------------
\subsubsection{Multipole expansion and the long-wavelength
               approximation}
\label{sec:multipole}
%--------------------------------------------------------------------

The spatial factor $e^{\iunit\mathbf{q}\cdot\mathbf{r}}$ in the
vector potential can be expanded in a Taylor series:
\begin{equation}
  e^{\iunit\mathbf{q}\cdot\mathbf{r}}
  = 1 + \iunit\mathbf{q}\cdot\mathbf{r}
  - \tfrac{1}{2}(\mathbf{q}\cdot\mathbf{r})^2
  + \cdots
  \label{eq:multipole}
\end{equation}
Each term corresponds to a class of multipole transitions with
distinct selection rules and characteristic coupling strengths.

The \textbf{zeroth-order term} (1) retains only the temporal
oscillation of $\Avec$, making it spatially uniform:
$\Avec(\mathbf{r},t) \to \Avec(t)$.  This is the
\textbf{long-wavelength} or \textbf{dipole approximation}.
The resulting $(e/mc)\Avec(t)\cdot\phat$ interaction drives
electric dipole (E1) transitions---the dominant contribution to
optical absorption.  The approximation is valid when the spatial
extent of the electronic system is much smaller than the photon
wavelength: $|\mathbf{q}||\mathbf{r}| \ll 1$.
For visible light ($\lambda \approx 400$--700~nm) interacting with
molecules (${\sim}1$~nm) or even nanoscale systems (${\sim}10$~nm),
the ratio $|\mathbf{q}||\mathbf{r}| \approx 10^{-2}$--$10^{-1}$ is
small and the approximation is excellent.

The \textbf{first-order term} ($\iunit\mathbf{q}\cdot\mathbf{r}$)
gives the electric quadrupole (E2) and magnetic dipole (M1)
contributions.  These are suppressed by a factor of
$|\mathbf{q}||\mathbf{r}| \approx v/c$ relative to E1 and are
important for forbidden transitions, circular dichroism, and X-ray
spectroscopy where $\lambda$ approaches atomic dimensions.

Higher-order terms give electric octupole, magnetic quadrupole,
etc., relevant only for very short wavelengths or very extended
systems.

The long-wavelength approximation (retaining only the zeroth-order
term) has been the standard in molecular quantum chemistry, including
in paper~1 of this series.\cite{jakowski-tddft-2025}
However, it has well-known limitations that become critical for
periodic systems, as discussed next.


%--------------------------------------------------------------------
\subsubsection{Failure of the length gauge for periodic systems}
\label{sec:lg-failure}
%--------------------------------------------------------------------

In the long-wavelength limit with spatially uniform $\Avec(t)$, two
equivalent representations of the light-matter interaction are
available, connected by the G\"{o}ppert-Mayer gauge transformation
(derived in Section~\ref{sec:GM-transform}):

\paragraph{Velocity gauge (VG).}
The interaction enters through the kinetic energy as
$(e/mc)\Avec(t)\cdot\phat$.  The perturbation operator involves
$\phat = -\iunit\hbar\nabla$, which is well-defined for any basis
set or boundary condition.

\paragraph{Length gauge (LG).}
A gauge transformation with $\chi(\mathbf{r},t)
= -\mathbf{r}\cdot\Avec(t)$ eliminates $\Avec$ from the kinetic
energy and transfers the interaction to a scalar potential
$\Phi = \Evec(t)\cdot\mathbf{r}$, giving the familiar dipole
interaction $-e\,\Evec(t)\cdot\hat{\mathbf{r}}$.  The perturbation
operator now involves the position operator $\hat{\mathbf{r}}$.

\medskip
\noindent
Note that these two gauges both operate within the Coulomb gauge
framework---they are different representations of the same physics
obtained by choosing different gauge functions, not different
physical theories.
For \emph{finite} molecular systems, both give identical results,
and the length gauge is often preferred because the dipole matrix
elements $\langle\psi_a|\hat{r}_\alpha|\psi_b\rangle$ have a
transparent physical interpretation as transition dipole moments.
This was the approach used in paper~1.\cite{jakowski-tddft-2025}

For \emph{periodic} systems under Born-von K\'{a}rm\'{a}n boundary
conditions, the length gauge fails catastrophically.
The position operator $\hat{\mathbf{r}}$ is fundamentally
incompatible with periodicity: $\mathbf{r}$ and $\mathbf{r}+\mathbf{L}$
must be physically equivalent, but $\hat{\mathbf{r}}$ distinguishes
them.
The scalar potential $\Evec\cdot\mathbf{r}$ is a linearly growing
ramp that produces a discontinuous ``sawtooth'' at the unit cell
boundary---it is not periodic and cannot be represented in a Bloch
basis.
The matrix elements
$\langle\psi_{n\mathbf{k}}|\hat{r}_\alpha|\psi_{m\mathbf{k}}\rangle$
between Bloch states are ill-defined (this is precisely the problem
resolved by the modern theory of polarization via Berry phases,
discussed in Section~\ref{sec:observables-berry}).

The velocity gauge completely avoids this problem: the momentum
operator $\phat$ is perfectly well-defined for Bloch states, and the
matrix elements
$\langle\psi_{n\mathbf{k}}|\hat{p}_\alpha|\psi_{m\mathbf{k}}\rangle$
exist without any difficulty.  This is why periodic-system RT-TDDFT
implementations, including the present one, use the velocity gauge.


%--------------------------------------------------------------------
\subsubsection{Scope and limitations of the velocity gauge with the
               dipole approximation}
\label{sec:vg-scope}
%--------------------------------------------------------------------

The velocity gauge resolves the periodic-boundary problem, but the
underlying long-wavelength approximation itself imposes a
fundamental limitation: it restricts the calculation to
\textbf{vertical transitions} with zero crystal momentum transfer
($\Delta\mathbf{k} = 0$).  This means that only the optical limit
of the dielectric function $\varepsilon(\mathbf{q}\to 0,\omega)$ is
accessible.

It is sometimes stated that the velocity gauge ``is not general,''
but this is imprecise.  The gauge transformation itself is
exact---gauge transformations are always unitary and introduce no
approximation.  What is limited is the \emph{dipole approximation}
that precedes it: by setting $e^{\iunit\mathbf{q}\cdot\mathbf{r}}
\to 1$, one discards all finite-$\mathbf{q}$ physics regardless of
which gauge is used.  The velocity gauge is the correct and rigorous
framework for $q = 0$ optical response in periodic systems; the
limitation lies in the $q = 0$ restriction, not in the gauge.

Many important spectroscopies and collective excitations involve
finite momentum transfer and lie beyond this approximation.
These include plasmon dispersion $\varepsilon(\mathbf{q},\omega)$,
electron energy-loss spectroscopy (EELS) where the beam transfers
finite $\mathbf{q}$ to the sample, inelastic X-ray scattering (IXS),
and magnon dispersion in magnetic systems.


%--------------------------------------------------------------------
\subsubsection{Beyond the dipole approximation: the full vector
               potential}
\label{sec:beyond-dipole}
%--------------------------------------------------------------------

A natural generalization is to retain the full spatial dependence of
the vector potential without making the dipole approximation:
\begin{equation}
  \hat{H}(t)
  = \frac{1}{2m}
    \!\left(\phat + \frac{e}{c}\Avec(\mathbf{r},t)\right)^{\!2}
  + V_\text{eff}[\rho](\mathbf{r},t)
  \label{eq:H-full-A}
\end{equation}
with $\Avec(\mathbf{r},t) = A_0\hat{\varepsilon}\,
e^{\iunit(\mathbf{q}\cdot\mathbf{r}-\omega t)} + \text{c.c.}$
retaining the full plane-wave spatial structure.  This is still the
Coulomb gauge---no new gauge choice is needed.  The spatial factor
$e^{\iunit\mathbf{q}\cdot\mathbf{r}}$ is automatically periodic (it
is a reciprocal-space function commensurate with the lattice for
appropriate $\mathbf{q}$), so the problematic $\hat{\mathbf{r}}$
operator never appears.  Transitions now satisfy crystal momentum
conservation $\mathbf{k}_f = \mathbf{k}_i + \mathbf{q}$
(mod~$\mathbf{G}$), allowing finite-$\mathbf{q}$ excitations.

This approach is more general than the dipole approximation because
the full exponential automatically includes all multipole orders.
For finite but small $|\mathbf{q}||\mathbf{r}|$, the leading
correction beyond E1 is the electric quadrupole / magnetic dipole
term (first order in $\mathbf{q}\cdot\mathbf{r}$).
For $|\mathbf{q}||\mathbf{r}| \sim 1$ (as in X-ray spectroscopy or
large nanostructures), all orders contribute and the full exponential
must be retained.

\paragraph{Important distinction regarding finite-$\mathbf{q}$
response.}
There is a second, independent route to finite-$\mathbf{q}$
excitations that does not involve the vector potential at all.
One can apply a scalar potential perturbation
\begin{equation}
  \delta V(\mathbf{r},t)
  = V_0\,e^{\iunit\mathbf{q}\cdot\mathbf{r}}\,\delta(t)
  \label{eq:density-wave-kick}
\end{equation}
directly to $V_\text{eff}$, with no gauge considerations.
This ``density-wave kick'' probes the density-density response
function $\chi(\mathbf{q},\omega)$, whereas the EM route probes the
current-current response $\sigma(\mathbf{q},\omega)$.
The two response functions are related but distinct:
$\sigma$ is a tensor ($3\times3$), $\chi$ is a scalar;
$\sigma$ naturally gives the transverse (optical) response,
$\chi$ gives the longitudinal response relevant to EELS.
For the longitudinal response at finite $\mathbf{q}$, the scalar
perturbation approach is more direct.
For transverse (optical) response including finite-$\mathbf{q}$
corrections (beyond dipole), the full
$\Avec(\mathbf{r},t)$ approach is the natural choice.

The present implementation uses the long-wavelength limit
(spatially uniform $\Avec(t)$) exclusively.
The general framework described here indicates the natural extension
path: keeping the full $e^{\iunit\mathbf{q}\cdot\mathbf{r}}$
spatial structure in $\Avec$ for beyond-dipole optical response, or
using scalar density-wave perturbations for longitudinal excitations
at arbitrary $\mathbf{q}$.
Such extensions would enable simulation of plasmon dispersion, EELS
spectra, inelastic X-ray scattering, and magnon dispersion from
first principles---all within the same RT-TDDFT framework.


%--------------------------------------------------------------------
\subsubsection{Observables: from dipole moment to current density to
               Berry phase polarization}
\label{sec:observables-berry}
%--------------------------------------------------------------------

The choice of gauge determines not only how the perturbation enters
the Hamiltonian but also which observable provides the most natural
route to spectral properties.
Three complementary approaches exist, corresponding to three
equivalent but operationally distinct ways of characterizing the
electronic response.

\paragraph{Length gauge---dipole moment.}
In the length gauge, the EM field couples through the scalar
potential $\Evec(t)\cdot\hat{\mathbf{r}}$, acting directly on the
charge density.  The natural observable is the time-dependent dipole
moment:
\begin{equation}
  \mu_\alpha(t)
  = -e\sum_n f_n
    \langle\psi_n(t)|\hat{r}_\alpha|\psi_n(t)\rangle,
  \label{eq:dipole-moment}
\end{equation}
where the sum runs over occupied KS orbitals with occupations $f_n$.
The Fourier transform of $\mu(t)$ after an impulsive perturbation
yields the frequency-dependent polarizability tensor $\alpha(\omega)$,
from which the photoabsorption cross section follows as
$\sigma_\text{abs}(\omega) = (4\pi\omega/c)\,
\text{Im}\,\alpha(\omega)$.
For finite molecular systems this is the standard RT-TDDFT
post-processing route, and it was used in
paper~1.\cite{jakowski-tddft-2025}

Higher-order multipole moments---the quadrupole
$Q_{\alpha\beta}(t) = \langle r_\alpha r_\beta\rangle$, the
octupole, etc.---can be extracted analogously and provide access to
E2, M1, and higher-order contributions to optical activity and
circular dichroism, though these are rarely needed for systems small
compared to the photon wavelength.

\paragraph{Velocity gauge---electronic current density.}
In the velocity gauge, the EM field couples through
$\Avec(t)\cdot\phat$, acting on the electronic momentum.
The natural observable is the time-dependent electronic current
density:
\begin{equation}
  \mathbf{j}(\mathbf{r},t)
  = \frac{1}{2}\sum_n f_n
    \left[\psi_n^*(\mathbf{r},t)\,\phat\,\psi_n(\mathbf{r},t)
          + \text{c.c.}\right]
  + \mathbf{j}_\text{NL}(\mathbf{r},t).
  \label{eq:current-density}
\end{equation}
The first term is the familiar paramagnetic current---the same
quantity that appears in probability current expressions in
introductory quantum mechanics.
The second term, $\mathbf{j}_\text{NL}$, arises from the nonlocal
pseudopotential: because the Kleinman-Bylander projectors do not
commute with the position operator, the continuity equation
$\partial\rho/\partial t + \nabla\cdot\mathbf{j} = 0$ is satisfied
only when the nonlocal contribution is included.
The derivation and explicit form of $\mathbf{j}_\text{NL}$ are
given in Section~\ref{sec:current}.
\myblue{[Verify cross-reference label for Sec.~2.5.]}

The macroscopic (cell-averaged) current
$\mathbf{J}(t) = (1/\Omega)\int_\text{cell}\mathbf{j}(\mathbf{r},t)
\,d^3r$
connects to optical properties through the frequency-dependent
optical conductivity:
\begin{equation}
  \sigma_{\alpha\beta}(\omega)
  = \frac{J_\alpha(\omega)}{E_\beta(\omega)},
  \label{eq:conductivity}
\end{equation}
and the dielectric function:
\begin{equation}
  \varepsilon_{\alpha\beta}(\omega)
  = \delta_{\alpha\beta}
  + \frac{4\pi\iunit}{\omega}\,\sigma_{\alpha\beta}(\omega),
  \label{eq:dielectric}
\end{equation}
from which absorption spectra, reflectivity, refractive index, and
energy-loss functions follow by standard relations.
This is the periodic-system analog of the molecular polarizability
route: where molecular RT-TDDFT extracts $\alpha(\omega)$ from
$\mu(t)$, periodic RT-TDDFT extracts $\varepsilon(\omega)$ from
$\mathbf{J}(t)$.

\paragraph{Berry phase---macroscopic polarization.}
The macroscopic polarization $\mathbf{P}$ of a periodic system
cannot be defined as the expectation value of any local
operator---this is the same fundamental difficulty with
$\hat{\mathbf{r}}$ discussed in Section~\ref{sec:lg-failure}.
The modern theory of polarization\cite{king-smith-vanderbilt-1993,%
resta-1994} resolves this by expressing $\mathbf{P}$ as a geometric
phase (Berry phase) of the occupied Bloch orbitals:
\begin{equation}
  P_\alpha
  = \frac{-e}{(2\pi)^3}\sum_n
    \int_\text{BZ}\!d\mathbf{k}\;
    \langle u_{n\mathbf{k}}|
    \iunit\frac{\partial}{\partial k_\alpha}
    |u_{n\mathbf{k}}\rangle,
  \label{eq:berry-phase-P}
\end{equation}
where $u_{n\mathbf{k}}(\mathbf{r})$ is the cell-periodic part of the
Bloch orbital.
This expression is defined modulo a quantum of polarization
$e\mathbf{R}/\Omega$ (where $\mathbf{R}$ is a lattice vector and
$\Omega$ is the cell volume), reflecting the intrinsic
multivaluedness of the geometric phase.

The Berry phase polarization is the periodic-system generalization
of the molecular dipole moment.
Just as $\mu = -e\langle\hat{\mathbf{r}}\rangle$ is a single
well-defined number for a localized molecule, $P^\text{Berry}$ is a
well-defined number (modulo the quantum) for an infinite periodic
crystal.
The Berry phase accomplishes for periodic systems what the position
operator does for molecules---but through a topological construction
in $k$-space rather than a real-space expectation value.

In RT-TDDFT, the time-evolved Bloch orbitals $\psi_{n\mathbf{k}}(t)$
at each time step define a time-dependent Berry phase, giving the
polarization $\mathbf{P}(t)$ directly without reference to the
position operator.
The connection to the current is exact:
\begin{equation}
  \mathbf{J}(t) = \frac{d\mathbf{P}}{dt},
  \label{eq:J-dPdt}
\end{equation}
so the polarization can equivalently be recovered by time-integration
of the current:
\begin{equation}
  \mathbf{P}(t) = \mathbf{P}(0) + \int_0^t\!\mathbf{J}(t')\,dt'.
  \label{eq:P-from-J}
\end{equation}
The two routes---direct Berry phase evaluation at each step versus
current integration---are formally equivalent but have different
numerical characteristics.
The Berry phase gives $\mathbf{P}(t)$ without accumulating
integration error, but requires dense $k$-point sampling (the
formula is a $k$-space integral) and careful tracking of the
polarization quantum for large-amplitude perturbations.
Current integration avoids the Berry phase calculation but
accumulates time-stepping error over long propagations.
In the linear response regime (small perturbations), both approaches
are reliable and serve as mutual consistency checks.
The Berry phase formalism and its implementation in the present code
are described in detail in Section~\ref{sec:berry-phase}.
\myblue{[Verify cross-reference label for Sec.~2.8.]}

\paragraph{Equivalence, complementarity, and the choice of
observable.}
For finite systems where $\hat{\mathbf{r}}$, $\phat$, and the Berry
phase are all well-defined, the three routes are exactly equivalent,
connected by the time derivative of the dipole moment:
\begin{equation}
  \frac{d\mu_\alpha}{dt}
  = -e\sum_n f_n
    \langle\psi_n(t)|\hat{v}_\alpha|\psi_n(t)\rangle
  = \Omega\,J_\alpha(t)
  = \Omega\,\frac{dP_\alpha}{dt}.
  \label{eq:mu-J-P-equiv}
\end{equation}
In frequency space this yields the fundamental relation
$\sigma(\omega) = \iunit\omega\,\alpha(\omega)$, linking the
polarizability to the conductivity, and the dielectric function
$\varepsilon(\omega) = 1 + 4\pi\alpha(\omega)/\Omega
= 1 + 4\pi\iunit\sigma(\omega)/\omega$.

For periodic systems, the dipole moment is undefined, leaving two
viable routes: the current density $\mathbf{J}(t)$, which gives
$\sigma(\omega)$ and $\varepsilon(\omega)$ directly via Fourier
transform; and the Berry phase $\mathbf{P}(t)$, which gives
$\alpha(\omega)$ and connects to the absolute polarization.
In the present work, we use the current density as the primary
observable (Sections~\ref{sec:current}--\myblue{[2.6]}) and the
Berry phase as a complementary diagnostic and for properties
where the absolute polarization is needed
(Section~\ref{sec:berry-phase}).

There is an appealing conceptual symmetry in the
gauge-observable pairing: in the length gauge, the field perturbs
the density (through $\Evec\cdot\hat{\mathbf{r}}$) and the response
is read from the density's moments; in the velocity gauge, the field
perturbs the current (through $\Avec\cdot\phat$) and the response is
read from the current.
The Berry phase stands apart from this pattern---it is a
gauge-invariant geometric property of the occupied manifold that can
be evaluated regardless of how the perturbation was applied,
depending only on the instantaneous state of the system, not on the
history of how it was driven.


%--------------------------------------------------------------------
\subsubsection{Summary and scope of the present work}
\label{sec:em-summary}
%--------------------------------------------------------------------

The general TDKS equation with full minimal coupling to
$\Avec(\mathbf{r},t)$ in the Coulomb gauge provides a unified
framework for modeling the interaction of electronic systems with
electromagnetic radiation.
The key approximations adopted in the present work and their
consequences are as follows.

The \textbf{Coulomb gauge} ($\nabla\cdot\Avec = 0$,
$\Phi_\text{rad} = 0$) carries all radiation physics in $\Avec$
alone, avoiding ambiguity between scalar and vector potential
contributions.
Both electric and magnetic fields are derivable from $\Avec$ via
standard Maxwell relations.

The \textbf{long-wavelength approximation}
($e^{\iunit\mathbf{q}\cdot\mathbf{r}} \to 1$) reduces
$\Avec(\mathbf{r},t)$ to a spatially uniform $\Avec(t)$,
restricting the calculation to vertical transitions
($\Delta\mathbf{k} = 0$).
This is equivalent to the electric dipole approximation and is valid
when the photon wavelength far exceeds the system size.

The \textbf{velocity gauge} represents the interaction as
$(e/mc)\Avec(t)\cdot\phat$ in the kinetic energy, avoiding the
ill-defined position operator that plagues the length gauge in
periodic systems.

The \textbf{current density} is the primary observable, giving the
optical conductivity $\sigma(\omega)$ and dielectric function
$\varepsilon(\omega)$ via Fourier transform.
The \textbf{Berry phase polarization} provides a complementary route
to the same spectral information and additionally gives the absolute
macroscopic polarization.

The present work operates within all three approximations: Coulomb
gauge, long-wavelength limit, velocity gauge.
This gives access to the full optical response $\varepsilon(\omega)$
of periodic systems with correct treatment of local field effects,
current contributions from nonlocal pseudopotentials, and
spin-collinear dynamics.
Extensions to finite-$\mathbf{q}$ response---via either the full
$\Avec(\mathbf{r},t)$ or scalar density-wave perturbations---and to
higher-multipole processes are natural next steps that the framework
accommodates without conceptual difficulty.
